\documentclass[9pt,twocolumn,a4paper]{article}
\usepackage[utf8]{inputenc}
\usepackage[T1]{fontenc}
\usepackage[nynorsk]{babel}
\usepackage{lmodern}
\usepackage{microtype}
\usepackage{amsmath,amssymb,amsthm,bm,bbm}
\usepackage[margin=1.35cm,top=1.4cm,bottom=1.4cm,columnsep=0.55cm]{geometry}
\usepackage{enumitem}
\setlist{nosep,leftmargin=1.4em}
\usepackage{titlesec}
\titlespacing*{\section}{0pt}{1.6ex plus 0.4ex minus 0.2ex}{0.8ex plus 0.2ex}
\titleformat{\section}{\normalfont\fontsize{10}{11}\bfseries}{\thesection.}{0.5em}{}

\newtheoremstyle{formstyle}{4pt}{4pt}{\itshape}{0pt}{\bfseries}{.}{0.4em}{}
\newtheoremstyle{defstyle}{4pt}{4pt}{\normalfont}{0pt}{\bfseries}{.}{0.4em}{}
\theoremstyle{defstyle}
\newtheorem{definisjon}{Definisjon}
\newtheorem{postulat}{Postulat}
\theoremstyle{formstyle}
\newtheorem{teorem}{Teorem}
\newtheorem{folgje}{Følgje}

\setlength{\abovedisplayskip}{4pt plus 1pt minus 1pt}
\setlength{\belowdisplayskip}{4pt plus 1pt minus 1pt}
\setlength{\parindent}{0.9em}
\setlength{\parskip}{1pt}

\title{\vspace{-1.2cm}\textbf{\large formlære}\\[0.2em]
\normalsize Ein allmenn teori om form, tilpassingslandskap og avgrensa agens}
\author{\small Iver Finne \\ \scriptsize Oslo, 2026}
\date{}

\begin{document}
\maketitle
\vspace{-0.8cm}

\begin{abstract}
\noindent Form er ei ordning av skilnader, realisert mellom moglege alternativ under samtidige trykk, gjennom agentar med avgrensa omhug og kapasitet. Her vert formgenerering, epistemisk konseptutvikling og seleksjon bundne saman i femtuppelet $(\mathbb{M}, \Pi, SG, \Lambda, \{A_i\})$. For eit endeleg formrom gjev maksimal relativ entropi til ei grunnfordeling, under ei grense for venta toleranseskalert avvik, ei Gibbs-fordeling. Ein lokal overgangskjerne kan utforske denne fordelinga når rom, vekter og estimat står faste, Hastings-korreksjonen sikrar detaljert balanse og dei reversible stega gjev samanheng og aperiodisitet. Kovariansrelasjonane mellom prioritering, estimert avvik og faktisk svikt følgjer då av eksponentialfamilien. Ei irreversibel realiseringshistorie krev ei eiga dynamisk skildring.
\end{abstract}

\section{Skilnad, struktur og substans}

\begin{definisjon}[Form, konfigurasjon og objekt]
Å skilje noko ut er å halde noko anna utanfor; det som står mellom dei, er ein \emph{skilnad}. Kvar side er det ho er, berre mot den andre. Skilnader i bestemde samanhengar utgjer ei \emph{ordning}, eller \emph{form}. Ein \emph{konfigurasjon} er ei bestemt ordning av bestemde delar; \emph{formverda} er totaliteten av realiserte konfigurasjonar. Ingen konfigurasjon er delvis realisert: anten er han tilfelle, eller han er det ikkje. Å forklare kvifor ein konfigurasjon vart realisert, er å vise kva vilkår og overgangar som førte til han, og å skilje mellom alternativ som var \emph{umoglege}, \emph{utilgjengelege}, \emph{uprøvde} og \emph{forkasta}.
\end{definisjon}

Ein konfigurasjonsdel som ikkje vert delt vidare på det valde analysenivået, er eit \emph{objekt}. Under kvart valt nivå held strukturen fram. Måten objekta heng saman på i konfigurasjonen, er \emph{strukturen} hans; \emph{forma} til konfigurasjonen er strukturen teken uavhengig av dei eigenskapane ved objekta analysen ser bort frå. Moglegheita eit objekt har for å inngå i konfigurasjonar på nivået, er \emph{valensen} til objektet. Når objektdeling, relasjonar og operasjonar er fastsette, avgrensar valensane eit moglegheitsrom på nivået. Ei anna oppdeling kan gjere andre delformer og operasjonar synlege.

\begin{postulat}[Det vidare moglege]
Det moglege er ikkje laga, men gjeve. Dette er eit ontologisk postulat, ikkje ei følgje av ei bestemt objektdeling. Under dette postulatet vert moglegheita avdekt og realiseringa laga; postulatet garanterer ikkje at eit fullstendig kart kan lagast på førehand.
\end{postulat}

\begin{definisjon}[Måling og oppleving]
Substans er den indre naturen til det som står i relasjonane, teke uavhengig av korleis ledda er ordna. Ei \emph{måling} er ein relasjon mellom det som vert målt og det som måler. Sidan fysikken uttrykkjer alt i målte storleikar og relasjonar mellom dei, seier fysikken kva det målte \emph{gjer}, ikkje kva leddet sjølv \emph{er}. Det som er kjent utan måling, ikkje gjennom ein ytre målerelasjon, men ved å vere det, er \emph{oppleving}. Dette skil tilgangsmåtar; det seier ikkje at oppleving manglar indre skilnader eller relasjonar. Kvar måling er gjeven som oppleving. Oppleving framstilt på nytt i seg sjølv er \emph{medvit}; oppleving føreset ikkje medvit.
\end{definisjon}

\begin{postulat}[Minimalitet, suffisiens og felt]
(i) Når ein kjend natur strekkjer til som indre natur for alle relasjonane modellen skal forklare, legg denne læra ingen annan substans til. (ii) Den relasjonelle strukturen i den fysiske skildringa kan vere ordning i oppleving utan at dei målte relasjonane må endrast. Dette suffisienspostulatet gjeld meir enn det einskilde subjektet kjenner innanfrå; det kan ikkje utleiast frå denne tilgangen åleine. (iii) Opplevinga er eitt felt; det som vert skilt ut i feltet, er mønster i det, ikkje bitar av det.
\end{postulat}

\begin{folgje}
Under suffisienspostulatet kan den fysiske strukturen vere ordna oppleving. Minimalitetspostulatet vel då oppleving som den eine substansen i denne læra. Slutninga er vilkårsbunden; ho er ikkje eit empirisk prov mot andre ontologiar som gjev dei same målbare relasjonane. Når ei avgrensing skil eit indre frå eit ytre, er det å kjenne henne innanfrå å vere henne; å kjenne henne utanfrå er å kjenne relasjonane ho står i. Ei avgrensing som held seg sjølv, har under individuasjonspostulatet tilgang innanfrå til si eiga avgrensing. Feltet som heilskap har inga utside og er ikkje eit slikt avgrensa sjølv. Materie er oppleving sett utanfrå, og form er korleis éin del av substansen tek seg ut for ein annan. Måten feltet oppfører seg regelmessig på, er naturlovene.
\end{folgje}

\section{Formrommet $\mathbb{M}$, delingane og formalgebraen}

\begin{definisjon}[Klasse og formrom]
Ei mengd konfigurasjonar som tillèt same operasjonar (delt \emph{affordansestruktur}), er ein \emph{klasse}. Mengda av alle moglege konfigurasjonar i klassen er \emph{formrommet} $\mathbb{M}$. I analysen vert rommet avgrensa ved klasse, nivå og operasjonar; eit nivåskifte kan krevje eit anna operativt rom. Kvar konfigurasjon er eitt punkt $x \in \mathbb{M}$; eit empirisk formrom er ein projeksjon av $\mathbb{M}$ langs eit valt sett målbare aksar $X_k$.
\end{definisjon}

\begin{definisjon}[Materiell og epistemisk deling]
Ved tidspunkt $t$ skil me mellom ei materiell og ei epistemisk deling av det oppgjevne formrommet $\mathbb{M}$. Skiljet inneber ikkje statistisk uavhengigheit:
\begin{enumerate}
\item \textbf{Materiell (ontologisk) deling:} Gjeven ved lovene i feltet brukte på materiala og operasjonane som er for hand ved $t$:
\begin{itemize}
\item \emph{Busette regionar} $B_t \subset \mathbb{M}$: posisjonar som er eller har vore realiserte (\emph{formhistoria}). Naboar er posisjonar skilde med eitt tillate steg.
\item \emph{Forbodne regionar} $V_t \subset \mathbb{M}$: logisk skildra posisjonar som er materielt utelukka under vilkåra ved $t$.
\item \emph{Opne regionar} $O_t = \mathbb{M} \setminus (B_t \cup V_t)$: ikkje-busette posisjonar som er materielt tillatne. Ein open region som kan nåast frå $B_t$ utan å krysse $V_t$, er \emph{tilgjengeleg}; elles \emph{utilgjengeleg} utan materialskifte. Dei opne posisjonane som grensar direkte til $B_t$, er \emph{det nærliggande moglege} $\partial B_t$.
\end{itemize}
\item \textbf{Epistemisk deling:} For ein agent $A$, eit kriterium $k$ og tidspunkt $t$, under oppgjevne prøvevilkår og tidshorisont, er $\Pi_{A,k,t} = (K_{A,k,t}, C_{A,k,t}, U_{A,k,t})$ ein partisjon av $\mathbb{M}$:
\begin{itemize}
\item \emph{Kunnskapsrommet} $K_{A,k,t}$: posisjonar der kriteriet er avgjort under prøvevilkåra; dei som oppfyller det, ligg i $K_{A,k,t}^{+}$.
\item \emph{Tankerommet} $C_{A,k,t}$: posisjonar som er tenkte, men der kriteriet ikkje er avgjort ($K_{A,k,t} \cap C_{A,k,t} = \emptyset$).
\item \emph{Det utenkjelege} $U_{A,k,t} = \mathbb{M} \setminus (K_{A,k,t} \cup C_{A,k,t})$: kan innanfrå berre nemnast som komplement. Når eit alternativ vert tenkt i høve til kriteriet, forlèt det denne mengda.
\end{itemize}
\end{enumerate}
\end{definisjon}

Der agent, prøvevilkår og tidspunkt står faste, skriv me $\Pi_k=(K_k,C_k,U_k)$ per omsyn og $\Pi=\{\Pi_k\}$ for den samla kunnskapsstanden. Omsyn som enno ikkje er skilde ut, høyrer til det akselause $U$. Når òg kriteriet står fast, skriv me $K$, $C$ og $U$. Same form kan vere kjend langs eitt omsyn og uavgjord eller utenkt langs eit anna. Ein forboden region kan vere tenkt ($V_t \cap C \neq \emptyset$), og ein open region kan vere utenkjeleg ($O_t \cap U \neq \emptyset$). Epistemisk rørsle har fire retningar: $C \to C$ (konseptekspansjon), $C \to K$ (avgjerd for kriteriet), $K \to C$ (opning av nye konjekturar) og $K \to K$ (oppløysingsauke). Snittet $F_{A,k,t} = C_{A,k,t} \cap \partial B_t$ er \emph{fronten}: tilgjengeleg og tenkt, men ikkje avgjord for kriteriet. Dette er ei operasjonell tilpassing av C--K-teorien til formposisjonar og kriterium.

\begin{definisjon}[Formalgebra og formgrammatikk]
Konfigurasjonar lukka under sum $+$, differanse $-$, produkt $\times$ og euklidiske eller affine transformasjonar $\theta$, med innleiringsrelasjonen $\leq$, utgjer ein \emph{formalgebra}. Der delane er liner, flater og lekamar, ser formalgebraen gjennom objektdelinga på nivået og finn delformer som inga førehandsdeling gav. Ein historisk \emph{formgrammatikk} er eit trippel $SG = (S, R, s_0)$ der $S \subseteq \mathbb{M}$, $R = \{a_r \to b_r\}$ er omskrivingsreglar, og $s_0 \in S$ er startforma:
\begin{align*}
\text{Utløysing:} &\quad \theta(a_r) \leq s \quad (\text{noverande geometri er einaste premiss}), \\
\text{Regelbruk:} &\quad s \to \big(s - \theta(a_r)\big) + \theta(b_r), \\
\text{Språk:} &\quad L(SG) = \{s \in S : s_0 \to^* s\}.
\end{align*}
Eitt regelbruk er eitt steg; $\partial B_t$ ligg eitt steg frå $B_t$. Ein grammatikkoperasjon kan gjere ei moglegheit tenkt ($C$); realisering og prøving gjennomfører $C \to K$ berre for det kriteriet resultatet avgjer under prøvevilkåra.
\end{definisjon}

\section{Seleksjonstrykk og tilpassingslandskapet $\Lambda$}

Affordansestrukturen som definerer ein klasse, er konstant innanfor klassen. Sidan ein konstant ikkje kan forklare ein variabel, seier det klassen er til for, ingenting om kvifor formene i han er ulike. Skilnadene krev \emph{seleksjonstrykk} som endrar sannsynsfordelinga over $\mathbb{M} \setminus V_t$ relativt til ei grunnfordeling $\mu(x)$.

\begin{postulat}[Additivitetsvilkåret og tilpassingslandskapet]
Dersom kvart verksamt trykk $j$ ved tid $t$ endrar vekta i $x$ med ein positiv faktor $f_j(x, t)$ og faktorane verkar multiplikativt til $\mu(x)\prod_j f_j(x, t)$, definerer $\lambda_j(x, t) = -\ln f_j(x, t)$ \emph{tilpassingslandskapet}:
\begin{equation}
\Lambda(x, t) = \sum_j \lambda_j(x, t).
\end{equation}
Der trykka samverkar ikkje-multiplikativt, må samspelsledda skiljast ut som eigne ledd i summen, eller faktorane definerast betinga av ei fastsett referanseordning. Indeksen $j$ omfattar då både einskildtrykk og dei særskilt oppgjevne samspelsledda; gradientoppdelinga følgjer den same faktor- og referanseordninga.
\end{postulat}

Verdien $\Lambda(x, t)$ er \emph{spenninga} i posisjonen $x$. Ved lik grunnvekt er lågare spenning meir sannsynleg; der ein gradient eksisterer, peikar $-\nabla\Lambda$ mot lågare samla spenning. Eit \emph{lokalt minimum} er ein posisjon utan betre tillatne nabosteg; ein \emph{dal} er attraktorregionen som ein gjeven navigasjonsmåte fører ned mot minimumet. Bratte dalveggar og etterfølgjande korrigerande rørsle \emph{kanaliserer} og fryser dimensjonane; slake dalar rommar fluktuasjon.

\begin{definisjon}[Stilart og landskapsdynamikk]
Formene si sverming kring eit felles tyngdepunkt i $\Lambda$ er ein \emph{stilart}. Stilen kan samle spor av mange samtidige trykk i éin etikett; som slik samlevariabel er han eit symptom på trykka. At stilperiode innanfor ein klasse med fast beregeometri predikerer formvariasjon betre enn materiale åleine, er ei hypotese som krev dei same målbare aksane og prøving på former modellane ikkje er tilpassa til.
\end{definisjon}

Kvar realisert form tek plass, bind stoff og endrar tyngda i $B_t$ gjennom verktøy, vane og forventing. Realiseringa endrar difor $\Lambda(x, t)$ lokalt: landskapet har minne (\emph{stiavhengigheit}), og all formgjeving er omformgjeving. Når gradvise trykkforskyvingar får ein dal til å forsvinne, gjer minimumet eit sprang (\emph{bifurkasjon}), som i formhistoria viser seg som brå skifte (radiasjon etterfølgd av konvergens) etter lange periodar med stillstand.

\section{Hinne, sjølv, omhug, lyskjegle og kapasitet}

\begin{definisjon}[Hinne og sjølv]
Ei \emph{hinne} (Markov-teppe) er ei mengd tilstandar som skil eit indre frå eit ytre slik at indre og ytre er vilkårsbestemt uavhengige når tilstandane i hinna er gjevne. Gjennom hinna går observasjonar inn og handlingar ut. Ei hinne er \emph{sjølvhalden} i den grad regulering innanfor henne motverkar forstyrringar som elles ville løyse henne opp. Eit \emph{sjølv} er eit mønster i opplevingsfeltet avgrensa av ei sjølvhalden hinne.
\end{definisjon}

\begin{postulat}[Individuasjon]
Den sjølvhaldne regulerande avgrensinga utanfrå er òg ei avgrensing innanfrå: sjølvet er ei avgrensa oppleving som berre når resten av feltet gjennom tilstandane på hinna. Persepsjon er difor eit grensesnitt, ikkje eit vindauge; kroppen er korleis den avgrensa opplevinga tek seg ut utanfrå.
\end{postulat}

\begin{definisjon}[Agent, spenning og horisontar]
Ein \emph{agent} er eit tilbakekopla system $A = (m, \pi, G, \varepsilon)$, der modellen $m$ knyter observasjonar og handlingar til venta følgjer, $G$ er målregionen, $\pi \in \mathcal{P}_A$ er politikken agenten tek i bruk, og $\varepsilon_k$ er toleransen langs omsyn $X_k$. Definisjonen er uavhengig av substrat; sjølve substratet (dei materielle vilkåra som gjev motstand, resonans og veto $V_t$ utan modell og politikk) er ein agent av nullte orden.
\end{definisjon}

Lat $e_k(t; \pi, d)$ vere den faktiske avstanden frå $G$ langs $X_k$ ved tid $t$ under politikken $\pi$ og forstyrringa $d$, og lat $D_k = D_{\mathrm{KL}}(p_{\mathrm{obs}} \parallel p_m)$ vere prediksjonsfeilen. Den \emph{indre spenninga} $\sigma_k(A)$ kjem frå tre kjelder som krev ulike svar:
\begin{enumerate}
\item \emph{Mål-avvik} $e_k$ (verda mot målet) $\to$ krev handling som endrar vilkåra.
\item \emph{Prediksjonsfeil} $D_k$ (sansinga mot forventinga) $\to$ krev oppdatering av modellen $m$.
\item \emph{Konflikt mellom omsyn} (mål langs éin akse mot mål langs ein annan) $\to$ krev kompromiss eller utviding av stegrepertoaret.
\end{enumerate}
Låg $D_k$ er ikkje låg $e_k$: ein agent kan føreseie samanbrotet presist utan å kunne hindre det. For eit sjølv er registrert mål-avvik over $\varepsilon_k$ kjent innanfrå som \emph{uro}, og under $\varepsilon_k$ som \emph{ro}. \emph{Haldshorisonten} $\tau_k(A)$ er kor lenge ein politikk $\pi \in \mathcal{P}_A$ kan halde $e_k(t; \pi, d) \leq \varepsilon_k$; \emph{prediksjonshorisonten} $\tau_k^{\mathrm{pred}}(A)$ er kor lenge $m$ føreseier følgjene med oppgjeven presisjon. Lyskjeglehorisonten $\tau_k^{\Gamma}(A)$ seier kor langt målet langs omsynet rekk fram i tid. Dei tre horisontane må prøvast kvar for seg, med same starttid og oppgjevne vilkår: kva agenten bryr seg om, kan føreseie og kan halde, treng ikkje ha same rekkjevidd.

\begin{definisjon}[Aksevis kapasitet og samla hald]
Lat $\kappa(A)$ vere aksane agenten har ein tilgjengeleg politikk som kan halde kvar for seg, og $\omega_k(A)>0$ tyde at han bryr seg om omsynet. Då er
\[
H_0(A)=\{X_k\in\kappa(A):\omega_k(A)>0\}.
\]
Eit samla hald krev éin politikk som held aksane samstundes under dei same forstyrringane $\mathcal D$ og i det same tidsrommet $T$:
\[
H(A;\pi)=\{X_k\in H_0(A):e_k(t;\pi,d)\leq\varepsilon_k
\text{ for alle }d\in\mathcal D,\ t\in T\}.
\]
For politikken agenten faktisk tek i bruk, skriv me $H(A)=H(A;\pi)$. To omsyn kan liggje i $H_0(A)$ kvar for seg utan å kunne haldast samstundes.
\end{definisjon}

For ein konsekvens som viser seg etter $\Delta_k$, er $\rho_k=\Delta_k/\tau_k^{\mathrm{pred}}(A)$ når prediksjonshorisonten er positiv. Ved $\rho_k\leq1$ kan følgja føreseiast med oppgjeven presisjon. Dersom $\Delta_k>\tau_k^\Gamma(A)$, fell ho utanfor målhorisonten. Med eit tilgjengeleg vern er dette \emph{likegyldigheit i tid}; utan både mål og kapasitet er det \emph{blindskap}. Dersom $\Delta_k\leq\tau_k^\Gamma(A)$, men $\Delta_k>\tau_k(A)$, er ho innanfor målet og utanfor haldet: \emph{avmakt i tid}. Når begge horisontane rekk, avgjer den nytta politikken kva som vert halde. Ved $\rho_k>1$ kan lokal tilbakekopling eller eit vern sett i verk no likevel halde omsynet når følgja kjem.
\end{document}
